% GENERATED FILE. Edit canonical lessons, then run npm run generate:books.
% canonical-source-hash: fnv1a64:5c6f35cf049d7d09
% canonical-lessons: AR-C483-buruda, AR-C483-barid, AR-C483-qaris, AR-C483-saqi, AR-C483-barada

\chapter{Cold}
\label{ch:ar-cold}
\hlchaptermodality{483}

\begin{chapteropening}
\textbf{By the end of this chapter:} \emph{I can say coldness, chilly, biting, frost and to get cold.}

Complete the last lesson of chapter 483: I can say coldness, chilly, biting, frost and to get cold.
\end{chapteropening}

\section[\texorpdfstring{\ar{برودة} (burūda)}{برودة (burūda)}]{\ar{برودة} (burūda) — coldness; a chill}

\label{lesson:AR-C483-buruda}



\begin{quote}
Before the new one: say the Arabic for to grow stronger, then the Arabic for danger.
\end{quote}

\subsection*{You'll want to know: \ar{برودة}}
\textbf{\ar{برودة}} — \emph{burūda} — \textquotedblleft{}coldness; a chill\textquotedblright{}.

\begin{quote}
\ar{شعرت} \ar{ببرودة} \ar{في} \ar{الغرفة}.
\end{quote}

\begin{quote}
\emph{shaʿartu bi-burūda fī al-ghurfa.}
\end{quote}

\begin{quote}
I felt a chill in the room.
\end{quote}

\begin{cousinweb}
\textbf{\ar{برودة}} (\emph{burūda}) is built on the native Arabic root \emph{b-r-d}, \textquotedblleft{}to be cold\textquotedblright{} (Lane). It is a verbal noun.
\end{cousinweb}

\begin{grammarlens}[title={What you've built}]
One more word for getting around.
\end{grammarlens}

\subsection*{Your turn}
\begin{itemize}
\raggedright
  \item \emph{Say it:} \emph{burūda}
  \item \emph{Say it:} \emph{burūda}, once more
  \item \emph{Say it:} \emph{shaʿartu bi-burūda fī al-ghurfa.} — I felt a chill in the room.
  \item \emph{Recall:} say the Arabic for to grow stronger, then the Arabic for danger, then say \emph{burūda} again
\end{itemize}

\begin{tcolorbox}[breakable,colback=teal!4,colframe=teal!35!black,title={Before you move on}]
What does \ar{برودة} mean? (\textquotedblleft{}Coldness; a chill\textquotedblright{}.) Say it once more.
\end{tcolorbox}
\section[\texorpdfstring{\ar{بارد} (bārid)}{بارد (bārid)}]{\ar{بارد} (bārid) — chilly; cold}

\label{lesson:AR-C483-barid}



\begin{quote}
Before the new one: say the Arabic for danger, then the Arabic for coldness.
\end{quote}

\subsection*{You'll want to know: \ar{بارد}}
\textbf{\ar{بارد}} — \emph{bārid} — \textquotedblleft{}chilly; cold\textquotedblright{}.

\begin{quote}
\ar{كان} \ar{الماء} \ar{باردا} \ar{جدا}.
\end{quote}

\begin{quote}
\emph{kāna al-māʾ bāridan jiddan.}
\end{quote}

\begin{quote}
The water was very cold.
\end{quote}

\begin{cousinweb}
\textbf{\ar{بارد}} (\emph{bārid}) is built on the native Arabic root \emph{b-r-d}, \textquotedblleft{}to be cold\textquotedblright{} (Lane). It is the active participle.
\end{cousinweb}

\begin{grammarlens}[title={What you've built}]
One more word for getting around.
\end{grammarlens}

\subsection*{Your turn}
\begin{itemize}
\raggedright
  \item \emph{Say it:} \emph{bārid}
  \item \emph{Say it:} \emph{bārid}, once more
  \item \emph{Say it:} \emph{kāna al-māʾ bāridan jiddan.} — The water was very cold.
  \item \emph{Recall:} say the Arabic for danger, then the Arabic for coldness, then say \emph{bārid} again
\end{itemize}

\begin{tcolorbox}[breakable,colback=teal!4,colframe=teal!35!black,title={Before you move on}]
What does \ar{بارد} mean? (\textquotedblleft{}Chilly; cold\textquotedblright{}.) Say it once more.
\end{tcolorbox}
\section[\texorpdfstring{\ar{قارس} (qāris)}{قارس (qāris)}]{\ar{قارس} (qāris) — biting (of cold); bitter}

\label{lesson:AR-C483-qaris}



\begin{quote}
Before the new one: say the Arabic for coldness, then the Arabic for chilly.
\end{quote}

\subsection*{You'll want to know: \ar{قارس}}
\textbf{\ar{قارس}} — \emph{qāris} — \textquotedblleft{}biting (of cold); bitter\textquotedblright{}.

\begin{quote}
\ar{كان} \ar{الشتاء} \ar{قارسا} \ar{في} \ar{الجبل}.
\end{quote}

\begin{quote}
\emph{kāna ash-shitāʾ qārisan fī al-jabal.}
\end{quote}

\begin{quote}
The winter was bitterly cold in the mountains.
\end{quote}

It goes with cold: \emph{bard qāris}, a biting cold.

\begin{cousinweb}
\textbf{\ar{قارس}} (\emph{qāris}) is built on the native Arabic root \emph{q-r-s}, \textquotedblleft{}to be intense (of cold); to freeze\textquotedblright{} (Lane). It is the active participle.
\end{cousinweb}

\begin{grammarlens}[title={What you've built}]
One more word for getting around.
\end{grammarlens}

\subsection*{Your turn}
\begin{itemize}
\raggedright
  \item \emph{Say it:} \emph{qāris}
  \item \emph{Say it:} \emph{qāris}, once more
  \item \emph{Say it:} \emph{kāna ash-shitāʾ qārisan fī al-jabal.} — The winter was bitterly cold in the mountains.
  \item \emph{Recall:} say the Arabic for coldness, then the Arabic for chilly, then say \emph{qāris} again
\end{itemize}

\begin{tcolorbox}[breakable,colback=teal!4,colframe=teal!35!black,title={Before you move on}]
What does \ar{قارس} mean? (\textquotedblleft{}Biting (of cold); bitter\textquotedblright{}.) Say it once more.
\end{tcolorbox}
\section[\texorpdfstring{\ar{صقيع} (ṣaqīʿ)}{صقيع (ṣaqīʿ)}]{\ar{صقيع} (ṣaqīʿ) — frost}

\label{lesson:AR-C483-saqi}



\begin{quote}
Before the new one: say the Arabic for chilly, then the Arabic for biting.
\end{quote}

\subsection*{You'll want to know: \ar{صقيع}}
\textbf{\ar{صقيع}} — \emph{ṣaqīʿ} — \textquotedblleft{}frost\textquotedblright{}.

\begin{quote}
\ar{كان} \ar{الصقيع} \ar{على} \ar{السيارة} \ar{في} \ar{الصباح}.
\end{quote}

\begin{quote}
\emph{kāna aṣ-ṣaqīʿ ʿalā as-sayyāra fī aṣ-ṣabāḥ.}
\end{quote}

\begin{quote}
There was frost on the car in the morning.
\end{quote}

\begin{cousinweb}
\textbf{\ar{صقيع}} (\emph{ṣaqīʿ}) is usually traced to the native Arabic root \emph{ṣ-q-ʿ}, \textquotedblleft{}hoarfrost\textquotedblright{}.
\end{cousinweb}

\begin{grammarlens}[title={What you've built}]
One more word for getting around.
\end{grammarlens}

\subsection*{Your turn}
\begin{itemize}
\raggedright
  \item \emph{Say it:} \emph{ṣaqīʿ}
  \item \emph{Say it:} \emph{ṣaqīʿ}, once more
  \item \emph{Say it:} \emph{kāna aṣ-ṣaqīʿ ʿalā as-sayyāra fī aṣ-ṣabāḥ.} — There was frost on the car in the morning.
  \item \emph{Recall:} say the Arabic for chilly, then the Arabic for biting, then say \emph{ṣaqīʿ} again
\end{itemize}

\begin{tcolorbox}[breakable,colback=teal!4,colframe=teal!35!black,title={Before you move on}]
What does \ar{صقيع} mean? (\textquotedblleft{}Frost\textquotedblright{}.) Say it once more.
\end{tcolorbox}
\section[\texorpdfstring{\ar{برد} (barada)}{برد (barada)}]{\ar{برد} (barada) — to get cold; to cool}

\label{lesson:AR-C483-barada}



\begin{quote}
Before the new one: say the Arabic for biting, then the Arabic for frost.
\end{quote}

\subsection*{You'll want to know: \ar{برد}}
\textbf{\ar{برد}} — \emph{barada} — \textquotedblleft{}to get cold; to cool\textquotedblright{}.

\begin{quote}
\ar{برد} \ar{الشاي} \ar{قبل} \ar{أن} \ar{أشربه}.
\end{quote}

\begin{quote}
\emph{barada ash-shāy qabla an ashrabahu.}
\end{quote}

\begin{quote}
The tea got cold before I drank it.
\end{quote}

Present: \ar{يبرد}, \emph{yabrud}.

The noun \emph{bard} is cold, and \emph{barad} is hail.

\begin{cousinweb}
\textbf{\ar{برد}} (\emph{barada}) is built on the native Arabic root \emph{b-r-d}, \textquotedblleft{}to be cold\textquotedblright{} (Lane). Lane gives \emph{barada} and \emph{baruda}, with the present \emph{yabrudu}.
\end{cousinweb}

\begin{grammarlens}[title={What you've built}]
One more word for getting around.
\end{grammarlens}

\subsection*{Your turn}
\begin{itemize}
\raggedright
  \item \emph{Say it:} \emph{barada}
  \item \emph{Say it:} \emph{barada}, once more
  \item \emph{Say it:} \emph{barada ash-shāy qabla an ashrabahu.} — The tea got cold before I drank it.
  \item \emph{Recall:} say the Arabic for biting, then the Arabic for frost, then say \emph{barada} again
\end{itemize}

\begin{tcolorbox}[breakable,colback=teal!4,colframe=teal!35!black,title={Before you move on}]
What does \ar{برد} mean? (\textquotedblleft{}To get cold; to cool\textquotedblright{}.) Say it once more.
\end{tcolorbox}
